\documentclass[10pt,a4paper]{amsart}
\usepackage[margin=19mm]{geometry}
\usepackage{amsmath,amssymb,mathtools}
\usepackage[scr=rsfs,cal=euler]{mathalfa}
\usepackage{xfrac}
\usepackage[unicode,hidelinks,bookmarks=false]{hyperref}
\newcommand{\ie}{i.e.\ }
\newcommand{\cf}{cf.\ }
\newcommand{\resp}{respectively\ }
\newcommand{\Subsection}{\subsection}
\providecommand{\nobreakdash}{\nobreak\hbox{-}}
\setlength{\emergencystretch}{2em}
\multlinegap0pt
\usepackage{float,adjustbox}
\usepackage{indentfirst,epic,xurl,graphics,needspace}
\setlength{\textwidth}{6.5in}
\setlength{\textheight}{9.7in}
\setlength{\evensidemargin}{0in}
\setlength{\oddsidemargin}{0in}
\setlength{\topmargin}{-.7in}
\setlength{\mathsurround}{.167em}
\newtheorem{thm}{Theorem}
\newtheorem{que}{Question}
\newtheorem{proj}{Project}
\newtheorem{theorem}{Theorem}[section]
\newtheorem{corollary}[theorem]{Corollary}
\newtheorem{proposition}[theorem]{Proposition}
\newtheorem{lemma}[theorem]{Lemma}
\newtheorem{hypothesis}[theorem]{Hypothesis}
\newtheorem{question}[theorem]{Question}
\newtheorem{example}[theorem]{Example}
\newtheorem{definition}[theorem]{Definition}
\def\bull{\vrule height .9ex width .8ex depth -.1ex }
\def\card#1{\vert #1 \vert}
\def\gpindex#1#2{\card {#1\colon #2}}
\def\irr#1{{\rm  Irr}(#1)}
\def\irri#1#2{{\rm Irr}_{(#1)} (#2)}
\def\cd#1{{\rm  cd}(#1)}
\def\cent#1#2{{\bf C}_{#1}(#2)}
\def\gpcen#1{{\bf Z} (#1)}
\def\ker#1{{\rm ker} (#1)}
\def\form#1#2#3{\langle\langle #1, #2 \rangle\rangle_{#3}}
\def\empform#1{\form {\cdot}{\cdot}{#1}}
\def\biform#1#2{\langle #1,#2 \rangle}
\def\empbiform{\biform {\cdot}{\cdot}}
\def\norm#1#2{{\rm N}_{#1} (#2)}
\def\NN{{\cal N}}
\def\MM{{\cal M}}
\def\TT{{\cal T}}
\def\UU{{\cal U}}
\def\QQ{{\cal Q}}
\def\PP{{\cal P}}
\def\BB{{\cal B}}
\newcommand{\stab}[2]{{#1}_{#2}}
\newcommand{\aut}[1]{{\rm Aut} (#1)}
\newcommand{\sym}[1]{{\rm Sym} (#1)}
\newcommand{\alt}[1]{{\rm Alt} (#1)}
\newcommand{\cn}[1]{{\rm cn} (#1)}
\newcommand{\psl}[2]{{\rm PSL}_{#1} (#2)}
\newcommand{\Bpi}[1]{{\rm B}_\pi (#1)}
\newcommand{\Ipi}[1]{{\rm I}_\pi (#1)}
\newcommand{\ZZ}{{\cal Z}}
\newcommand{\nl}[1]{{\rm nlIrr} (#1)}
\newcommand{\GG}{{\cal G}}
\newcommand{\HH}{{\cal H}}
\begin{document}
\title[Some answers regarding factorizations of finite groups]{Some answers regarding factorizations of finite groups}
\author{Ryan McCulloch}
\date{}
\maketitle
\noindent\emph{Source and licence.} Some answers regarding factorizations of finite groups. Version arXiv:2607.20569v3 (CC BY 4.0). This material is distributed under the Creative Commons Attribution 4.0 International licence, http://creativecommons.org/licenses/by/4.0/, as recorded in the arXiv OAI licence metadata for this exact version and shown on the abstract page https://arxiv.org/abs/2607.20569v3. Full rights notice: licenses/arxiv-2607.20569-CC-BY-4.0.txt.\par\smallskip
\noindent\textit{Editorial scope.} Counted unit: the original \texttt{theorem} environment \texttt{thm2} at source lines 149--167 together with its complete proof at lines 169--202; the delivered document keeps source lines 113--203 so that every referenced label and every citation resolves inside it. Native editable source is the paper's own concatenated TeX; the AMS wrapper is editorial formatting only. No publisher class, upstream script or unrelated artwork is redistributed.
\begin{lemma}\label{lem: first_last}
Let $A_1 \cdots A_k$ be a factorization of a finite group $G$. Then
\begin{enumerate}
\item $|A_1|$ divides the order of the subgroup of $G$ generated by the set ${A_1}^{-1} A_1 = \{ g^{-1}h \,\, | \,\, g,h \in A_1 \}$, which can also be described as generated by any one of the subsets $g^{-1} A_1$ $(g \in A_1)$.  Moreover, that order is also the order of the subgroup generated by $A_1{A_1}^{-1} = \{ gh^{-1} \,\, | \,\, g,h \in A_1 \}$, equivalently, by any one of the subsets $A_1 g^{-1}$ $(g \in A_1)$.
\item Similarly, $|A_k|$ divides the order of the subgroup of $G$ generated by the set ${A_k}^{-1} A_k$, equivalently, by any of the subsets $g^{-1} A_k$ $(g \in A_k)$, and that order is also the order of the subgroup generated by $A_k{A_k}^{-1}$, equivalently, by any of the subsets $A_k g^{-1}$ $(g \in A_k)$.
\end{enumerate}
\end{lemma}

We also need the following lemma which appears in \cite{berg, sands}. 

\begin{lemma}\label{lem: e}
If $G = A_1 \cdots A_k$ is a factorization of a group $G$, then for all $g,h \in G$, $$(gA_1) \cdot A_2 \cdots A_{k-1} \cdot (A_k h)$$ is also a factorization of G.

Hence if for some positive integers $n_1, \dots, n_k$, $G$ has a $k$-fold factorization with $|A_i| = n_i$ $(i = 1, \dots, k)$, it has such a factorization in which $A_1$ and $A_k$ both contain the identity element $e$.
\end{lemma}

We are now ready to present the counterexample.  Let $k \geq 3$ and let $G$ be the Frobenius group of order $2^{k-1}(2^{k-1}-1)$ with kernel $N$ an elementary abelian $2$-group of order $2^{k-1}$ and complement $A$ cyclic of order $2^{k-1}-1$.  Viewing $N$ as the additive group of the field $\mathbb{F}_{2^{k-1}}$ and viewing $A$ as generated by a primitive element in the multiplicative group of the field, we have $A$ acting on $N$ via field multiplication.  Two key properties of this group that we need are that the action of $A$ on $N$ is transitive on the nonidentity elements of $N$, and that the elements of $G$ of even order are precisely the nonidentity elements of $N$, all of which have order $2$.

When $k=3$ the example is the alternating group $A_4$ and as we will see, it is the same example given by Bergman in \cite{berg}.  Our arguments mirror those for the $k=3$ case in \cite{berg}.

\begin{theorem}\label{thm1}
Let $k \geq 3$ and let $G$ be the Frobenius group of order $2^{k-1}(2^{k-1}-1)$ as described above.  Then $G$ possesses no $k$-fold factorization for the factorization $(2,\dots,2,2^{k-1}-1,2)$ of $|G|$.
\end{theorem}

\begin{proof}
Suppose $G = A_1 \cdots A_k$ is a factorization with $|A_i|=2$ for all $1 \leq i \leq k-2$, $|A_{k-1}| = 2^{k-1} - 1$, and $|A_k| = 2$.  By Lemma $\ref{lem: e}$ we can assume without loss of generality that $A_1$ and $A_k$ have the forms $\{e,a\}$ and $\{e,b\}$ respectively.  By Lemma \ref{lem: first_last}, $|A_1|$ divides $|\langle {A_1}^{-1}A_1 \rangle| = |\langle a \rangle|$ and $|A_k|$ divides $|\langle {A_k}^{-1}A_k \rangle| = |\langle b \rangle|$.  The only even order elements of $G$ are the elements of $N$ which are all order $2$, hence we conclude that $A_1$ and $A_k$ are in fact subgroups of $N$ (which may or may not be distinct).

Write $S = A_2 \cdots A_{k-1}$.  Note that for each $g \in S$ we have $A_1gA_k \subseteq NgN = Ng = gN$ as $N$ is a normal subgroup of $G$.  And so in order for $G = A_1 \cdot S \cdot A_k$ to hold, we must have $S$ containing representatives from every coset of $N$ in $G$.  As the action of the complement $A$ is transitive on the nonidentity elements of $N$, there exists $g \in S$ such that $g$ conjugates $A_1$ to $A_k$.

Hence $A_1 g A_k = g A_k A_k$.  But the multiplication map $A_k \times A_k \rightarrow A_k$ is not one-to-one; and it follows that the multiplication map $A_1 \times \cdots \times A_k \rightarrow G$ is not one-to-one, contradicting $G = A_1 \cdots A_k$ being a factorization.
\end{proof}

\section{A Sufficient Condition for $k$-Factorizations to Exist}

In a recent note by Bergman \cite{berg2}, a sufficient condition is given for $2$-factorizations of a finite group to exist.  Bergman asks whether a similar approach can be used to give a sufficient condition for $k$-factorizations to exist for larger $k$.  Our theorem below is a direct generalization of Bergman's Theorem 1.1 in \cite{berg2}, and the argument mirrors his argument.  The $k=2$ case yields Bergman's result.


\begin{theorem}\label{thm2}
Let $G$ be a finite group, and
%
\begin{equation}\begin{minipage}[c]{35pc}\label{d.G0Gn}
$\{e\}\,=\,G_0\ <\ G_1\ <\ \dots\ <\ G_n\,=\,G$ \ $(n\geq 0)$
\end{minipage}\end{equation}
%
a chain of subgroups of $G$.

Then for any factorization $|G| = a_1 a_2 \cdots a_k$ with $a_1, \dots, a_k$ positive integers, a {\em sufficient} condition for there
to exist subsets $A_1,\dots,A_k \subseteq G$ of cardinalities $|A_i| = a_i$ for all $i$ such that $$G = A_1 \cdots A_k$$
is that there exist subsets $s_1 , \dots, s_k$ of $\{1,\dots,n\}$, $s_i \cap s_j = \emptyset$ for $i \neq j$ and $\bigcup_{i \in \{1,\dots,n\}} s_i = \{1,\dots,n\}$ such that 
%
\begin{equation*}\begin{minipage}[c]{35pc}
for each $j \in \{1,\dots,n\}$, if $j \in s_t$, then $s_a \cap \{1,\dots,j\} = \emptyset$ holds either for all $a < t$ or for all $a > t$,
\end{minipage}\end{equation*}
%
and such that for each $t \in \{1,\dots,k\}$, $\prod_{i\in s_t}\,|G_i:\,G_{i-1}| = a_t$.
\end{theorem}

\begin{proof}
Given a finite group $G$, a chain~\eqref{d.G0Gn}, and such subsets $s_1, \dots s_k$, we shall show by induction on $j$ that for each $j \leq n$, 
%
\begin{equation}\begin{minipage}[c]{35pc}\label{d.A1jA2j}
there exists a factorization $G_j = A_{1,j} \cdots A_{k,j}$, such that\\
 for each $t \in \{1,\dots,k\}$, $|A_{t,j}| = \Pi_{i \in s_t \cap \{1,\dots,j\}} | G_i : G_{i-1}|$.
\end{minipage}\end{equation}

This is trivially true for $j=0$.   (In that case the numerical products shown are over the empty set, hence
their values are all $1 = |\{e\}|$, and $G_0 = \{e\} = \{e\}\cdots\{e\}$.)

So let $j < n$, and assume inductively that we have a factorization $G_j = A_{1,j} \cdots A_{k,j}$ as in~\eqref{d.A1jA2j}.  

Now $j+1 \in s_t$ for some unique $t$.  We know that $s_a \cap \{1,\dots,j+1\} = \emptyset$ holds either for all $a < t$ or for all $a > t$.  Without loss of generality, suppose 
%
\begin{equation*}\begin{minipage}[c]{35pc}
$s_a \cap \{1,\dots,j+1\} = \emptyset$ for all $a < t$.
\end{minipage}\end{equation*}
%
Thus $|A_{a,j}| = 1$ for all $a < t$.  Write $S = A_{1,j} \cdots A_{t-1,j}$.  So 
\begin{equation}\begin{minipage}[c]{35pc}\label{d.i_in_s1}
$G_j = S A_{t,j} \cdots A_{k,j}$, and by Lemma \ref{lem: e}, we may assume $S = \{ e \}$.
\end{minipage}\end{equation}

Let $C_{j+1}$ be a full set of representatives of the left cosets of $G_j$ in $G_{j+1}$, and let
$A_{t,j+1} = C_{j+1} A_{t,j}$ and let $A_{i,j+1} = A_{i,j}$ for all $i \neq t$.  

Given that $A_{1,j} \dots A_{k,j}$have the cardinalities indicated in~\eqref{d.A1jA2j}, we see that the new sets $A_{1,j+1} \dots A_{k,j+1}$ will
have cardinalities as in the $j\,{+}\,1$ case of~\eqref{d.A1jA2j}.  Finally, assuming~\eqref{d.i_in_s1}, we obtain $$A_{1,j+1} \cdots A_{k,j+1} = C_{j+1} A_{t,j} \cdots A_{k,j} = C_{j+1} G_j = G_{j+1}.$$

So we have proved the $j + 1$ case of~\eqref{d.A1jA2j};
so by induction,~\eqref{d.A1jA2j} holds for all $j \leq n$. The case $j = n$ is
our desired conclusion.
\end{proof}


\begin{thebibliography}{99}
\bibitem[berg]{berg} G. M. Bergman, A note on factorizations of finite groups, {\it J. Iran. Math. Soc.} {\bf 1} (2020), 157--161.
\bibitem[berg2]{berg2} G. M. Bergman, An observation on factorizations of finite groups, unpublished note (last revised 7 August, 2026). [\url{https://math.berkeley.edu/~gbergman/papers/unpub/re_gp_factzn.pdf}]
\bibitem[sands]{sands} S. Szab\'o and A. D. Sands, Factoring Groups Into Subsets, Lecture Notes in Pure and Applied Mathematics, 257, CRC Press, Boca Raton, FL, 2009.
\end{thebibliography}
\end{document}
